\chapter*{Glossaire}
\addcontentsline{toc}{chapter}{Glossaire}

\begin{description}[leftmargin=3cm, labelwidth=2.5cm, font=\bfseries]

    \item[AES-256-GCM] \textit{Advanced Encryption Standard -- Galois/Counter Mode}.
    Algorithme de chiffrement symétrique opérant sur des blocs de 128 bits avec une clé
    de 256 bits. Le mode GCM assure simultanément la confidentialité et l'intégrité des
    données via un tag d'authentification.

    \item[API REST] \textit{Application Programming Interface -- Representational State Transfer}.
    Interface de programmation exposant des ressources via des URLs et des méthodes HTTP
    standard (GET, POST, PATCH, DELETE).

    \item[CI/CD] \textit{Continuous Integration / Continuous Delivery}.
    Pratique d'intégration continue consistant à automatiser les tests et
    le déploiement à chaque modification du code source.

    \item[CRUD] \textit{Create, Read, Update, Delete}.
    Ensemble des quatre opérations de base sur une ressource persistante.

    \item[DRF] \textit{Django REST Framework}.
    Bibliothèque Python construite sur Django facilitant le développement d'API REST :
    sérialisation, authentification, permissions, pagination.

    \item[Docker] Plateforme de conteneurisation permettant d'empaqueter une application
    et ses dépendances dans des conteneurs isolés, assurant la reproductibilité
    entre environnements (développement, production).

    \item[HTTPS] \textit{HyperText Transfer Protocol Secure}.
    Version chiffrée du protocole HTTP utilisant TLS/SSL pour sécuriser
    les échanges entre client et serveur.

    \item[JWT] \textit{JSON Web Token}.
    Standard ouvert (RFC 7519) définissant un format compact pour transmettre
    des informations de manière sécurisée entre parties sous forme de token signé
    (HMAC-SHA256 ou RSA).

    \item[Nginx] Serveur web et proxy inverse haute performance utilisé ici pour
    terminer le chiffrement TLS, rediriger HTTP vers HTTPS, et router les requêtes
    vers le frontend React et le backend Django.

    \item[ORM] \textit{Object-Relational Mapper}.
    Couche d'abstraction traduisant les objets Python en requêtes SQL.
    Django ORM élimine les requêtes SQL brutes et prévient les injections SQL.

    \item[RBAC] \textit{Role-Based Access Control}.
    Modèle de contrôle d'accès attribuant des droits aux utilisateurs via
    leurs rôles (ADMIN, MANAGER, ANALYST, VIEWER) plutôt qu'individuellement.

    \item[UUID] \textit{Universally Unique Identifier}.
    Identifiant de 128 bits généré aléatoirement (RFC 4122), utilisé comme clé primaire
    pour empêcher l'énumération des ressources (\texttt{/offres/1}, \texttt{/offres/2}).

    \item[XSS] \textit{Cross-Site Scripting}.
    Vulnérabilité web permettant à un attaquant d'injecter du code JavaScript malveillant
    dans les pages vues par d'autres utilisateurs. Stockage du token JWT en mémoire React
    (non \texttt{localStorage}) constitue une protection contre ce vecteur.

\end{description}

\newpage
